\begin{figure}[H]
\centering
\begin{tikzpicture}
\begin{axis}[
  width=15cm,
  height=9cm,
  xlabel={\rl{موج بارومتر عرب}},
  ylabel={\rl{درصد اعتماد}},
  symbolic x coords={I (06),II (11),III (13),IV (16),V (19),VII (22)},
  xtick=data,
  xticklabel style={font=\small},
  ymin=0, ymax=80,
  legend pos=outer north east,
  legend style={font=\tiny, cells={anchor=west}},
  grid=major,
  grid style={dashed, gray!25},
  title={\rl{\textbf{روند اعتماد به نهادها در عراق -- بارومتر عرب}}},
  title style={font=\small},
]

\addplot[color=olive, mark=*, very thick] coordinates {
  ({I (06)},55) ({II (11)},62) ({III (13)},58) ({IV (16)},65) ({V (19)},70) ({VII (22)},72)
};
\addlegendentry{\rl{ارتش}}

\addplot[color=blue, mark=square*, very thick] coordinates {
  ({I (06)},45) ({II (11)},32) ({III (13)},18) ({IV (16)},15) ({V (19)},10) ({VII (22)},15)
};
\addlegendentry{\rl{دولت}}

\addplot[color=red, mark=triangle*, very thick] coordinates {
  ({I (06)},28) ({II (11)},18) ({III (13)},10) ({IV (16)},8) ({V (19)},5) ({VII (22)},8)
};
\addlegendentry{\rl{احزاب سیاسی}}

\addplot[color=purple, mark=diamond*, very thick, dashed] coordinates {
  ({I (06)},60) ({II (11)},52) ({III (13)},45) ({IV (16)},40) ({V (19)},35) ({VII (22)},32)
};
\addlegendentry{\rl{رهبران مذهبی}}

\end{axis}
\end{tikzpicture}
\caption{روند اعتماد به نهادها -- بارومتر عرب (۲۰۰۶--۲۰۲۲)}
\label{fig:ab-trust-trend}
\end{figure}

\subsubsection{مهم‌ترین چالش‌ها از نظر شهروندان}

\textbf{متن دقیق پرسش:}
\begin{quote}
\lr{``What is the most important challenge facing your country?''}
\end{quote}

\begin{table}[H]
\centering
\caption{مهم‌ترین چالش -- بارومتر عرب}
\label{tab:ab-challenges}
\begin{tabularx}{\textwidth}{>{\bfseries\small}r *{6}{>{\centering\arraybackslash\small}X}}
\toprule
\textbf{چالش} & \textbf{I} & \textbf{II} & \textbf{III} & \textbf{IV} & \textbf{V} & \textbf{VII} \\
\midrule
امنیت/تروریسم & ۴۸٪ & ۲۵٪ & ۳۵٪ & ۴۰٪ & ۱۰٪ & ۸٪ \\
فساد مالی/اداری & ۸٪ & ۲۵٪ & ۲۸٪ & ۲۵٪ & ۴۵٪ & ۳۸٪ \\
بیکاری & ۱۵٪ & ۲۰٪ & ۱۵٪ & ۱۵٪ & ۲۰٪ & ۲۵٪ \\
وضعیت اقتصادی & ۱۲٪ & ۱۵٪ & ۱۰٪ & ۱۰٪ & ۱۵٪ & ۱۸٪ \\
خدمات عمومی & ۱۰٪ & ۸٪ & ۵٪ & ۵٪ & ۵٪ & ۵٪ \\
مداخله‌ی خارجی & ۵٪ & ۵٪ & ۵٪ & ۳٪ & ۳٪ & ۴٪ \\
تنش فرقه‌ای & ۲٪ & ۲٪ & ۲٪ & ۲٪ & ۲٪ & ۲٪ \\
\bottomrule
\end{tabularx}
\end{table}

\infobox{boxblue}{تحلیل: چرخش از امنیت به فساد}{
یکی از مهم‌ترین تحولات افکار عمومی عراق، چرخش اولویت از \textbf{امنیت} (۲۰۰۶) 
به \textbf{فساد} (۲۰۱۹) است. این چرخش نشان می‌دهد:
\begin{itemize}[rightmargin=0.5cm]
\item با کاهش نسبی خشونت، شهروندان فرصت یافتند به مشکلات ساختاری بیندیشند
\item فساد سیستماتیک نه‌تنها کاهش نیافته بلکه نهادینه‌تر شده
\item جنبش تشرین ۲۰۱۹ بر همین بستر نارضایتی از فساد شکل گرفت
\item مردم دریافته‌اند بسیاری از مشکلات امنیتی ریشه در فساد دارد
\end{itemize}
}

\subsubsection{ارزیابی عملکرد دولت}

\textbf{متن دقیق پرسش:}
\begin{quote}
\lr{``How would you evaluate the current government's performance in the following areas: job creation, controlling prices, fighting corruption, narrowing the gap between rich and poor, providing adequate health care, providing adequate education?''}
\end{quote}

\begin{table}[H]
\centering
\caption{ارزیابی عملکرد دولت (درصد «خوب/بسیار خوب») -- بارومتر عرب}
\label{tab:ab-govt-performance}
\small
\begin{tabularx}{\textwidth}{>{\bfseries\small}r *{4}{>{\centering\arraybackslash\small}X}}
\toprule
\textbf{حوزه} & \textbf{III (۱۳)} & \textbf{IV (۱۶)} & \textbf{V (۱۹)} & \textbf{VII (۲۲)} \\
\midrule
ایجاد شغل & ۸٪ & ۵٪ & ۵٪ & ۸٪ \\
کنترل قیمت‌ها & ۱۲٪ & ۸٪ & ۸٪ & ۱۰٪ \\
مبارزه با فساد & ۱۰٪ & ۵٪ & ۳٪ & ۸٪ \\
کاهش شکاف فقیر/غنی & ۸٪ & ۵٪ & ۳٪ & ۵٪ \\
خدمات بهداشتی & ۱۸٪ & ۱۲٪ & ۱۰٪ & ۱۵٪ \\
خدمات آموزشی & ۲۰٪ & ۱۵٪ & ۱۲٪ & ۱۵٪ \\
تأمین برق & ۱۵٪ & ۱۰٪ & ۱۲٪ & ۱۵٪ \\
تأمین آب & ۱۸٪ & ۱۲٪ & ۱۵٪ & ۱۸٪ \\
\bottomrule
\end{tabularx}
\end{table}


% ============================================================
% بخش ۵: نظرسنجی ارزش‌های جهانی (WVS)
% ============================================================
\section{نظرسنجی ارزش‌های جهانی (\lr{World Values Survey})}

\subsection{روش‌شناسی}

\begin{table}[H]
\centering
\caption{مشخصات \lr{WVS} در عراق}
\label{tab:wvs-method}
\begin{tabularx}{\textwidth}{>{\bfseries}r X}
\toprule
\textbf{ویژگی} & \textbf{شرح} \\
\midrule
مدیریت جهانی & دانشگاه‌های استکهلم و میشیگان \\
موج‌ها در عراق & موج ۵ (۲۰۰۴) و موج ۶ (۲۰۱۳) \\
حجم نمونه & \ptho{۲}{۳۲۵} (۲۰۰۴) و \ptho{۱}{۲۰۰} (۲۰۱۳) \\
ویژگی & عمیق‌ترین نظرسنجی ارزشی -- ۲۵۰+ سؤال \\
دسترسی & داده‌ها عمومی: \lr{worldvaluessurvey.org} \\
\bottomrule
\end{tabularx}
\end{table}

\subsection{ارزش‌های سنتی در برابر سکولار--عقلانی}

\begin{table}[H]
\centering
\caption{شاخص‌های ارزشی کلیدی \lr{WVS} -- عراق}
\label{tab:wvs-values}
\small
\begin{tabularx}{\textwidth}{>{\bfseries\small}r >{\small}X >{\centering\arraybackslash\small}p{1.5cm} >{\centering\arraybackslash\small}p{1.5cm}}
\toprule
\textbf{پرسش/شاخص} & \textbf{شرح} & \textbf{۲۰۰۴} & \textbf{۲۰۱۳} \\
\midrule
\multicolumn{4}{r}{\textbf{مذهب و سکولاریزاسیون}} \\
مذهب بسیار مهم است & \lr{``How important is religion in your life?''} & ۹۵٪ & ۹۲٪ \\
به خدا ایمان دارم & \lr{``Do you believe in God?''} & ۹۹٪ & ۹۹٪ \\
رهبران مذهبی بر سیاست تأثیر بگذارند & \lr{``Religious leaders should influence government''} & ۵۵٪ & ۴۵٪ \\
\midrule
\multicolumn{4}{r}{\textbf{دموکراسی و حکمرانی}} \\
دموکراسی مهم است & \lr{``How important is it to live in a democratic country?'' (7-10 از ۱۰)} & ۷۵٪ & ۶۵٪ \\
رهبر قوی بدون انتخابات & \lr{``Having a strong leader who doesn't bother with parliament''} (خوب) & ۳۵٪ & ۴۲٪ \\
حکومت نظامی & \lr{``Having the army rule''} (خوب) & ۲۲٪ & ۲۸٪ \\
حکومت کارشناسان & \lr{``Having experts make decisions''} (خوب) & ۶۸٪ & ۷۲٪ \\
\midrule
\multicolumn{4}{r}{\textbf{جنسیت و نقش زنان}} \\
مردان رهبران سیاسی بهتری هستند & \lr{``Men make better political leaders than women''} & ۸۲٪ & ۷۵٪ \\
تحصیلات دانشگاهی برای پسران مهم‌تر & \lr{``University education more important for boys''} & ۵۵٪ & ۴۵٪ \\
زنان باید حجاب داشته باشند & -- & ۸۰٪ & ۷۵٪ \\
\midrule
\multicolumn{4}{r}{\textbf{تحمل و پذیرش دیگری}} \\
نمی‌خواهم همسایه‌ام باشد: نژاد دیگر & -- & ۲۵٪ & ۲۰٪ \\
نمی‌خواهم همسایه‌ام باشد: مذهب دیگر & -- & ۲۸٪ & ۳۲٪ \\
نمی‌خواهم همسایه‌ام باشد: خارجی & -- & ۳۵٪ & ۴۰٪ \\
\bottomrule
\end{tabularx}
\end{table}

\infobox{boxyellow}{تحلیل \lr{WVS}: عراق در نقشه‌ی فرهنگی جهان}{
بر اساس نقشه‌ی فرهنگی اینگلهارت--ولزل (\lr{Inglehart-Welzel Cultural Map}):
\begin{itemize}[rightmargin=0.5cm]
\item عراق در ربع \textbf{ارزش‌های سنتی + ارزش‌های بقا} قرار دارد
\item در مقایسه با ۲۰۰۴، در ۲۰۱۳ حرکت اندکی به‌سمت ارزش‌های سکولار مشاهده شد
\item تمایل به «رهبر قوی» و «حکومت نظامی» \textbf{افزایش} یافته -- نشانه‌ی سرخوردگی از دموکراسی
\item حمایت از نقش سیاسی مذهب \textbf{کاهش} یافته -- نشانه‌ی سرخوردگی از اسلام سیاسی
\item تحمل «دیگری» (مذهب/نژاد/ملیت دیگر) \textbf{کاهش} یافته -- اثر جنگ فرقه‌ای
\end{itemize}
}


% ============================================================
% بخش ۶: نظرسنجی‌های Pew Research Center
% ============================================================
\section{نظرسنجی‌های \lr{Pew Research Center} (۲۰۰۳--۲۰۱۵)}

\subsection{روش‌شناسی}

\begin{table}[H]
\centering
\caption{مشخصات نظرسنجی‌های \lr{Pew} در عراق}
\label{tab:pew-method}
\begin{tabularx}{\textwidth}{>{\bfseries}r X}
\toprule
\textbf{ویژگی} & \textbf{شرح} \\
\midrule
مرکز & \lr{Pew Research Center} (واشنگتن) \\
بخش & \lr{Global Attitudes Project} \\
حجم نمونه & \ptho{۱}{۰۰۰} تا \ptho{۱}{۵۰۰} \\
سال‌ها & ۲۰۰۳ (محدود)، ۲۰۰۴، ۲۰۰۵، ۲۰۰۶، ۲۰۰۷، ۲۰۰۸، ۲۰۰۹، ۲۰۱۲، ۲۰۱۵ \\
\bottomrule
\end{tabularx}
\end{table}

\subsection{نگرش به ایالات متحده}

\textbf{متن دقیق پرسش:}
\begin{quote}
\lr{``Do you have a favorable or unfavorable opinion of the United States?''}
\end{quote}

\begin{table}[H]
\centering
\caption{نگرش عراقی‌ها به ایالات متحده -- \lr{Pew}}
\label{tab:pew-us-favorability}
\begin{tabularx}{\textwidth}{>{\bfseries}r *{7}{>{\centering\arraybackslash}X}}
\toprule
 & \textbf{۲۰۰۳} & \textbf{۲۰۰۴} & \textbf{۲۰۰۵} & \textbf{۲۰۰۷} & \textbf{۲۰۰۹} & \textbf{۲۰۱۲} & \textbf{۲۰۱۵} \\
\midrule
\multicolumn{8}{r}{\textbf{نگرش مثبت به آمریکا}} \\
کل & -- & ۴۰٪ & ۳۲٪ & ۲۲٪ & ۳۰٪ & ۳۲٪ & ۲۸٪ \\
شیعه & -- & ۴۵٪ & ۳۸٪ & ۲۵٪ & ۳۵٪ & ۳۵٪ & ۳۰٪ \\
سنی & -- & ۲۰٪ & ۱۲٪ & ۵٪ & ۱۲٪ & ۱۵٪ & ۱۲٪ \\
کرد & -- & ۶۵٪ & ۶۰٪ & ۵۵٪ & ۵۸٪ & ۵۵٪ & ۵۰٪ \\
\midrule
\multicolumn{8}{r}{\textbf{نگرش منفی به آمریکا}} \\
کل & -- & ۵۵٪ & ۶۲٪ & ۷۲٪ & ۶۴٪ & ۶۰٪ & ۶۵٪ \\
\midrule
\multicolumn{8}{r}{\textbf{مقایسه‌ی منطقه‌ای -- نگرش مثبت}} \\
اردن & ۱٪ & ۵٪ & ۲۱٪ & ۲۰٪ & ۲۵٪ & ۱۲٪ & ۱۴٪ \\
ترکیه & ۱۵٪ & ۳۰٪ & ۲۳٪ & ۹٪ & ۱۴٪ & ۱۵٪ & ۱۸٪ \\
لبنان & ۲۷٪ & -- & ۴۲٪ & ۴۷٪ & ۵۵٪ & ۴۸٪ & ۳۸٪ \\
\bottomrule
\end{tabularx}
\end{table}

\subsection{نگرش به دموکراسی}

\textbf{متن دقیق پرسش:}
\begin{quote}
\lr{``Some people believe that democracy is a Western way of doing things that would not work here. Others feel that democracy can work here. Which is closer to your view?''}

\textit{برخی معتقدند دموکراسی شیوه‌ی غربی است که اینجا جواب نمی‌دهد. 
برخی دیگر معتقدند دموکراسی اینجا هم می‌تواند کار کند. 
کدام به نظر شما نزدیک‌تر است؟}
\end{quote}

\begin{table}[H]
\centering
\caption{«آیا دموکراسی در اینجا کار می‌کند؟» -- \lr{Pew}}
\label{tab:pew-democracy-works}
\begin{tabularx}{\textwidth}{>{\bfseries}r *{5}{>{\centering\arraybackslash}X}}
\toprule
 & \textbf{۲۰۰۴} & \textbf{۲۰۰۵} & \textbf{۲۰۰۷} & \textbf{۲۰۰۹} & \textbf{۲۰۱۲} \\
\midrule
دموکراسی کار می‌کند & ۶۰٪ & ۵۵٪ & ۴۸٪ & ۵۵٪ & ۵۲٪ \\
شیوه‌ی غربی/کار نمی‌کند & ۳۰٪ & ۳۵٪ & ۴۰٪ & ۳۵٪ & ۳۸٪ \\
نمی‌دانم & ۱۰٪ & ۱۰٪ & ۱۲٪ & ۱۰٪ & ۱۰٪ \\
\midrule
\multicolumn{6}{r}{\textbf{«کار می‌کند» -- تفکیک قومی--مذهبی}} \\
شیعه & ۶۵٪ & ۶۰٪ & ۵۰٪ & ۶۰٪ & ۵۵٪ \\
سنی & ۴۰٪ & ۳۵٪ & ۳۰٪ & ۳۵٪ & ۳۵٪ \\
کرد & ۷۵٪ & ۷۰٪ & ۶۸٪ & ۷۰٪ & ۶۸٪ \\
\bottomrule
\end{tabularx}
\end{table}

\subsection{نگرش به تنش شیعه--سنی}

\textbf{متن دقیق پرسش:}
\begin{quote}
\lr{``Is there a struggle between modernizers and fundamentalists in our country?''}

\lr{``Are you very concerned, somewhat concerned, not too concerned, or not at all concerned about tensions between Sunni and Shia Muslims?''}
\end{quote}

\begin{table}[H]
\centering
\caption{نگرانی از تنش شیعه--سنی -- \lr{Pew}}
\label{tab:pew-sectarian}
\begin{tabularx}{\textwidth}{>{\bfseries}r *{4}{>{\centering\arraybackslash}X}}
\toprule
 & \textbf{۲۰۰۶} & \textbf{۲۰۰۷} & \textbf{۲۰۰۹} & \textbf{۲۰۱۲} \\
\midrule
\multicolumn{5}{r}{\textbf{نگرانی از تنش فرقه‌ای (مجموع «زیاد» و «نسبتاً»)}} \\
کل & ۸۸٪ & ۹۲٪ & ۸۰٪ & ۷۵٪ \\
شیعه & ۸۵٪ & ۹۰٪ & ۷۸٪ & ۷۲٪ \\
سنی & ۹۲٪ & ۹۵٪ & ۸۵٪ & ۸۰٪ \\
\midrule
\multicolumn{5}{r}{\textbf{وجود مبارزه بین مدرن‌گرایان و بنیادگرایان}} \\
بله & ۶۵٪ & ۷۰٪ & ۶۰٪ & ۵۵٪ \\
\bottomrule
\end{tabularx}
\end{table}


% ============================================================
% بخش ۷: نظرسنجی‌های IRI و NDI
% ============================================================
\section{نظرسنجی‌های \lr{IRI} و \lr{NDI} (۲۰۰۳--۲۰۱۳)}

\subsection{\lr{International Republican Institute (IRI)}}

\begin{table}[H]
\centering
\caption{مشخصات نظرسنجی‌های \lr{IRI}}
\label{tab:iri-method}
\begin{tabularx}{\textwidth}{>{\bfseries}r X}
\toprule
\textbf{ویژگی} & \textbf{شرح} \\
\midrule
سفارش‌دهنده & \lr{International Republican Institute} \\
هدف & ارزیابی فرایند دموکراسی‌سازی و عملکرد دولت \\
حجم نمونه & \ptho{۳}{۰۰۰} تا \ptho{۳}{۵۰۰} (بزرگ‌ترین نمونه‌ها) \\
موج‌ها & ۱۲ موج بین ۲۰۰۳ تا ۲۰۱۳ \\
ویژگی & پرسش‌های بسیار عملیاتی درباره‌ی خدمات دولتی \\
\bottomrule
\end{tabularx}
\end{table}

\subsubsection{ارزیابی خدمات عمومی}

\textbf{متن دقیق پرسش:}
\begin{quote}
\lr{``How would you rate the availability of the following services in your area: electricity, clean water, sewage, health care, education, employment opportunities? (Excellent, Good, Fair, Poor, Very Poor)''}
\end{quote}

\begin{table}[H]
\centering
\caption{ارزیابی خدمات عمومی (درصد «خوب/عالی») -- \lr{IRI}}
\label{tab:iri-services}
\small
\begin{tabularx}{\textwidth}{>{\bfseries\small}r *{6}{>{\centering\arraybackslash\small}X}}
\toprule
\textbf{خدمت} & \textbf{۲۰۰۴} & \textbf{۲۰۰۶} & \textbf{۲۰۰۸} & \textbf{۲۰۱۰} & \textbf{۲۰۱۱} & \textbf{۲۰۱۳} \\
\midrule
برق & ۱۵٪ & ۱۲٪ & ۱۸٪ & ۲۰٪ & ۲۲٪ & ۲۰٪ \\
آب آشامیدنی & ۳۵٪ & ۳۰٪ & ۳۵٪ & ۳۸٪ & ۴۰٪ & ۳۸٪ \\
فاضلاب & ۲۰٪ & ۱۸٪ & ۲۲٪ & ۲۵٪ & ۲۵٪ & ۲۲٪ \\
بهداشت و درمان & ۲۲٪ & ۱۸٪ & ۲۵٪ & ۲۸٪ & ۳۰٪ & ۲۵٪ \\
آموزش & ۳۰٪ & ۲۵٪ & ۳۵٪ & ۳۸٪ & ۳۸٪ & ۳۵٪ \\
فرصت‌های شغلی & ۱۰٪ & ۸٪ & ۱۲٪ & ۱۵٪ & ۱۵٪ & ۱۲٪ \\
\midrule
\multicolumn{7}{r}{\textbf{برق -- تفکیک استانی (۲۰۱۱)}} \\
بغداد & & & & & ۱۵٪ & \\
بصره & & & & & ۱۲٪ & \\
نجف & & & & & ۱۸٪ & \\
انبار & & & & & ۱۰٪ & \\
اربیل & & & & & ۵۵٪ & \\
سلیمانیه & & & & & ۵۰٪ & \\
\bottomrule
\end{tabularx}
\end{table}

\infobox{boxred}{یافته‌ی کلیدی \lr{IRI}: شکست خدمات عمومی}{
حتی ده سال پس از حمله‌ی ۲۰۰۳:
\begin{itemize}[rightmargin=0.5cm]
\item فقط ۲۰٪ از عراقی‌ها برق را «خوب» ارزیابی کردند
\item فقط ۱۲٪ فرصت‌های شغلی را مناسب دانستند
\item تفاوت فاحش بین کردستان (۵۵٪ رضایت از برق) و بقیه‌ی عراق (۱۰--۱۸٪)
\item این ارقام پس از صرف بیش از ۶۰ میلیارد دلار بودجه‌ی بازسازی به‌دست آمده
\item شکست خدمات عمومی یکی از ریشه‌های اصلی بی‌اعتمادی به حکومت است
\end{itemize}
}

\subsubsection{ارزیابی عملکرد مقامات منتخب}

\textbf{متن دقیق پرسش:}
\begin{quote}
\lr{``How would you rate the job performance of the following: Prime Minister, Provincial Governor, Provincial Council, Local Mayor/Council?''}
\end{quote}

\begin{table}[H]
\centering
\caption{ارزیابی عملکرد مقامات (درصد «خوب/عالی») -- \lr{IRI}}
\label{tab:iri-officials}
\small
\begin{tabularx}{\textwidth}{>{\bfseries\small}r *{6}{>{\centering\arraybackslash\small}X}}
\toprule
\textbf{مقام} & \textbf{۲۰۰۴} & \textbf{۲۰۰۶} & \textbf{۲۰۰۸} & \textbf{۲۰۱۰} & \textbf{۲۰۱۱} & \textbf{۲۰۱۳} \\
\midrule
نخست‌وزیر & ۵۰٪ & ۴۲٪ & ۴۵٪ & ۴۸٪ & ۳۵٪ & ۲۸٪ \\
استاندار & ۳۵٪ & ۳۰٪ & ۳۵٪ & ۳۸٪ & ۳۲٪ & ۲۸٪ \\
شورای استان & ۲۸٪ & ۲۲٪ & ۲۸٪ & ۳۰٪ & ۲۵٪ & ۲۰٪ \\
شهردار/شورای محلی & ۳۲٪ & ۲۵٪ & ۳۲٪ & ۳۵٪ & ۲۸٪ & ۲۵٪ \\
نماینده‌ی پارلمان & ۲۵٪ & ۱۸٪ & ۲۲٪ & ۲۵٪ & ۱۸٪ & ۱۵٪ \\
\midrule
\multicolumn{7}{r}{\textbf{نخست‌وزیر -- تفکیک قومی--مذهبی (۲۰۱۱: مالکی)}} \\
شیعه & & & & & ۴۵٪ & \\
سنی & & & & & ۱۲٪ & \\
کرد & & & & & ۳۵٪ & \\
\bottomrule
\end{tabularx}
\end{table}

\subsubsection{مشارکت سیاسی و انتخاباتی}

\textbf{متن دقیق پرسش‌ها:}
\begin{quote}
\lr{``Did you vote in the last election? Do you plan to vote in the next election?''}

\lr{``What is the most important factor in choosing a candidate?''}

\lr{``Do you believe your vote makes a difference?''}
\end{quote}

\begin{table}[H]
\centering
\caption{مشارکت سیاسی و رفتار انتخاباتی -- \lr{IRI}}
\label{tab:iri-participation}
\small
\begin{tabularx}{\textwidth}{>{\bfseries\small}r *{5}{>{\centering\arraybackslash\small}X}}
\toprule
\textbf{پرسش} & \textbf{۲۰۰۵} & \textbf{۲۰۰۸} & \textbf{۲۰۱۰} & \textbf{۲۰۱۱} & \textbf{۲۰۱۳} \\
\midrule
\multicolumn{6}{r}{\textbf{مشارکت انتخاباتی}} \\
در انتخابات قبلی رأی دادم & ۷۲٪ & ۶۸٪ & ۶۵٪ & ۵۸٪ & ۵۵٪ \\
قصد دارم در انتخابات بعدی رأی دهم & ۸۰٪ & ۷۵٪ & ۷۰٪ & ۶۵٪ & ۶۰٪ \\
رأی من تأثیر دارد & ۶۵٪ & ۵۵٪ & ۴۸٪ & ۴۲٪ & ۳۸٪ \\
\midrule
\multicolumn{6}{r}{\textbf{مهم‌ترین معیار انتخاب نامزد}} \\
برنامه و عملکرد & ۳۰٪ & ۲۵٪ & ۲۸٪ & ۳۰٪ & ۳۲٪ \\
تعلق فرقه‌ای/قومی & ۲۵٪ & ۳۰٪ & ۲۵٪ & ۲۲٪ & ۲۰٪ \\
تعلق حزبی & ۱۵٪ & ۱۵٪ & ۱۲٪ & ۱۰٪ & ۱۰٪ \\
صداقت و پاکدستی & ۲۰٪ & ۲۲٪ & ۲۵٪ & ۲۸٪ & ۳۰٪ \\
توصیه‌ی مرجع دینی & ۱۰٪ & ۸٪ & ۱۰٪ & ۱۰٪ & ۸٪ \\
\bottomrule
\end{tabularx}
\end{table}

\infobox{boxgreen}{روند مثبت: کاهش رأی فرقه‌ای}{
یکی از معدود روندهای مثبت: سهم «تعلق فرقه‌ای/قومی» به‌عنوان معیار 
انتخاب نامزد از ۳۰٪ (۲۰۰۸) به ۲۰٪ (۲۰۱۳) کاهش یافت. 
در مقابل، «صداقت و پاکدستی» از ۲۲٪ به ۳۰٪ افزایش یافت. 
این روند در انتخابات ۲۰۱۸ و ۲۰۲۱ و به‌ویژه در جنبش تشرین ادامه یافت.
}


\subsection{\lr{National Democratic Institute (NDI)}}

\begin{table}[H]
\centering
\caption{مشخصات نظرسنجی‌های \lr{NDI}}
\label{tab:ndi-method}
\begin{tabularx}{\textwidth}{>{\bfseries}r X}
\toprule
\textbf{ویژگی} & \textbf{شرح} \\
\midrule
سفارش‌دهنده & \lr{National Democratic Institute for International Affairs} \\
هدف & ارزیابی فرایند انتخابات و مشارکت مدنی \\
مجری & \lr{Greenberg Quinlan Rosner Research} \\
حجم نمونه & \ptho{۲}{۰۰۰} تا \ptho{۲}{۴۰۰} \\
موج‌ها & ۸ موج بین ۲۰۰۴ تا ۲۰۰۹ \\
\bottomrule
\end{tabularx}
\end{table}

\subsubsection{پرسش‌های کلیدی \lr{NDI}}

\textbf{پرسش ۱: آمادگی برای زندگی مسالمت‌آمیز}
\begin{quote}
\lr{``Can Sunnis, Shia, and Kurds live together peacefully in one country, or is the division too deep?''}
\end{quote}

\begin{table}[H]
\centering
\caption{«آیا شیعه، سنی و کرد می‌توانند مسالمت‌آمیز زندگی کنند؟» -- \lr{NDI}}
\label{tab:ndi-coexistence}
\begin{tabularx}{\textwidth}{>{\bfseries}r *{4}{>{\centering\arraybackslash}X}}
\toprule
 & \textbf{۲۰۰۴} & \textbf{۲۰۰۶} & \textbf{۲۰۰۸} & \textbf{۲۰۰۹} \\
\midrule
\multicolumn{5}{r}{\textbf{«بله، می‌توانند»}} \\
کل & ۷۵٪ & ۵۸٪ & ۵۵٪ & ۶۵٪ \\
شیعه & ۸۰٪ & ۶۰٪ & ۵۸٪ & ۶۸٪ \\
سنی & ۷۰٪ & ۵۲٪ & ۴۸٪ & ۵۸٪ \\
کرد & ۶۵٪ & ۵۰٪ & ۵۵٪ & ۶۲٪ \\
\midrule
\multicolumn{5}{r}{\textbf{«خیر، شکاف بسیار عمیق است»}} \\
کل & ۲۰٪ & ۳۵٪ & ۳۸٪ & ۲۸٪ \\
\bottomrule
\end{tabularx}
\end{table}

\textbf{پرسش ۲: ساختار مطلوب حکومت}
\begin{quote}
\lr{``Which kind of government do you think would be best for Iraq: a strong central government, a federal system, or separate states?''}
\end{quote}

\begin{table}[H]
\centering
\caption{ساختار حکومتی مطلوب -- \lr{NDI}}
\label{tab:ndi-government-structure}
\begin{tabularx}{\textwidth}{>{\bfseries}r *{4}{>{\centering\arraybackslash}X}}
\toprule
 & \textbf{۲۰۰۴} & \textbf{۲۰۰۶} & \textbf{۲۰۰۸} & \textbf{۲۰۰۹} \\
\midrule
\multicolumn{5}{r}{\textbf{کل}} \\
دولت مرکزی قوی & ۶۰٪ & ۵۵٪ & ۴۸٪ & ۵۲٪ \\
نظام فدرالی & ۲۵٪ & ۲۸٪ & ۳۲٪ & ۳۰٪ \\
دولت‌های جداگانه/تجزیه & ۱۰٪ & ۱۲٪ & ۱۵٪ & ۱۲٪ \\
نمی‌دانم & ۵٪ & ۵٪ & ۵٪ & ۶٪ \\
\midrule
\multicolumn{5}{r}{\textbf{«دولت مرکزی قوی» -- تفکیک}} \\
شیعه & ۵۵٪ & ۵۲٪ & ۴۵٪ & ۵۰٪ \\
سنی & ۸۵٪ & ۸۰٪ & ۷۵٪ & ۷۸٪ \\
کرد & ۲۰٪ & ۱۵٪ & ۱۲٪ & ۱۵٪ \\
\midrule
\multicolumn{5}{r}{\textbf{«دولت‌های جداگانه» -- تفکیک}} \\
شیعه & ۸٪ & ۱۰٪ & ۱۲٪ & ۱۰٪ \\
سنی & ۲٪ & ۳٪ & ۵٪ & ۴٪ \\
کرد & ۳۵٪ & ۴۰٪ & ۴۵٪ & ۳۸٪ \\
\bottomrule
\end{tabularx}
\end{table}

\infobox{boxyellow}{تحلیل: شکاف ساختاری عمیق}{
\begin{itemize}[rightmargin=0.5cm]
\item \textbf{سنی‌ها} شدیداً خواهان دولت مرکزی قوی (۷۵--۸۵٪) -- 
ترس از تهمیش در نظام فدرالی
\item \textbf{کردها} خواهان فدرالیسم یا حتی جدایی (۳۵--۴۵٪ خواهان تجزیه) -- 
تجربه‌ی خودمختاری موفق
\item \textbf{شیعیان} در میانه -- تمایل به دولت مرکزی اما با حفظ اکثریت سیاسی خود
\item این شکاف ساختاری، ریشه‌ی بسیاری از بن‌بست‌های سیاسی عراق است
\end{itemize}
}


% ============================================================
% بخش ۸: نظرسنجی‌های Oxford Research International
% ============================================================
\section{نظرسنجی‌های \lr{Oxford Research International} (۲۰۰۳--۲۰۰۴)}

\subsection{اهمیت تاریخی}
نظرسنجی‌های \lr{Oxford Research International (ORI)} اولین نظرسنجی‌های بزرگ‌مقیاس 
در عراق پس از سقوط صدام بودند و تصویری منحصربه‌فرد از لحظه‌ی گذار ارائه می‌دهند.

\begin{table}[H]
\centering
\caption{مشخصات نظرسنجی‌های \lr{ORI}}
\label{tab:ori-method}
\begin{tabularx}{\textwidth}{>{\bfseries}r X}
\toprule
\textbf{ویژگی} & \textbf{شرح} \\
\midrule
مجری & \lr{Oxford Research International} (آکسفورد) \\
سفارش‌دهنده & \lr{BBC, ABC News, NHK, ARD} \\
موج‌ها & اکتبر ۲۰۰۳، فوریه ۲۰۰۴، نوامبر ۲۰۰۴ \\
حجم نمونه & \ptho{۲}{۵۰۰}+ در هر موج \\
ویژگی خاص & اولین نظرسنجی‌ها پس از سقوط رژیم \\
\bottomrule
\end{tabularx}
\end{table}

\subsection{نتایج کلیدی موج اول (اکتبر ۲۰۰۳)}

\textbf{فقط ۶ ماه پس از سقوط صدام:}

\begin{table}[H]
\centering
\caption{نتایج کلیدی اولین نظرسنجی پس از سقوط (اکتبر ۲۰۰۳)}
\label{tab:ori-first}
\begin{tabularx}{\textwidth}{>{\bfseries\small}r X >{\centering\arraybackslash}p{2cm}}
\toprule
\textbf{پرسش} & \textbf{متن اصلی} & \textbf{پاسخ مثبت} \\
\midrule
شرایط زندگی بهتر شده & \lr{``Life is better now than before the war''} & ۵۸٪ \\
\addlinespace
از سرنگونی صدام خوشحالم & \lr{``I'm glad Saddam is gone''} & ۶۲٪ \\
\addlinespace
خوش‌بین به آینده & \lr{``Things will be better in a year''} & ۶۷٪ \\
\addlinespace
دموکراسی می‌خواهم & \lr{``I want democracy for Iraq''} & ۷۲٪ \\
\addlinespace
نیروهای ائتلاف باید بمانند & \lr{``Coalition forces should stay until stable''} & ۵۰٪ \\
\addlinespace
احساس آزادی بیشتر & \lr{``I feel freer now''} & ۶۰٪ \\
\addlinespace
اوضاع امنیتی خوب & \lr{``Security in my area is good''} & ۵۵٪ \\
\addlinespace
فساد بزرگ‌ترین مشکل & \lr{``Corruption is the biggest problem''} & ۸٪ \\
\addlinespace
بیکاری بزرگ‌ترین مشکل & \lr{``Unemployment is the biggest problem''} & ۲۸٪ \\
\addlinespace
امنیت بزرگ‌ترین مشکل & \lr{``Security is the biggest problem''} & ۳۵٪ \\
\bottomrule
\end{tabularx}
\end{table}

\infobox{boxblue}{مقایسه‌ی ۲۰۰۳ با ۲۰۱۹: ارزیابی شانزده‌ساله}{
\begin{table}[H]
\centering
\small
\begin{tabularx}{\textwidth}{>{\bfseries\small}r >{\centering\arraybackslash\small}X >{\centering\arraybackslash\small}X >{\centering\arraybackslash\small}X}
\toprule
\textbf{شاخص} & \textbf{۲۰۰۳ (\lr{ORI})} & \textbf{۲۰۱۹ (\lr{AB})} & \textbf{تغییر} \\
\midrule
خوش‌بینی به آینده & ۶۷٪ & ۲۰٪ & $\downarrow$ ۴۷ واحد \\
اعتماد به دولت & ۵۰٪ & ۱۰٪ & $\downarrow$ ۴۰ واحد \\
فساد مهم‌ترین مشکل & ۸٪ & ۴۵٪ & $\uparrow$ ۳۷ واحد \\
دموکراسی مطلوب & ۷۲٪ & ۵۵٪ & $\downarrow$ ۱۷ واحد \\
احساس آزادی & ۶۰٪ & ۴۵٪ & $\downarrow$ ۱۵ واحد \\
\bottomrule
\end{tabularx}
\end{table}
\vspace{4pt}
این مقایسه تصویری دراماتیک از \textbf{فرسایش امید} طی شانزده سال ارائه می‌دهد. 
خوش‌بینی اولیه‌ی ۲۰۰۳ که محصول سقوط دیکتاتوری و وعده‌ی آینده‌ی بهتر بود، 
در مواجهه با واقعیت خشونت، فساد و ناکارآمدی فروریخت.
}


% ============================================================
% بخش ۹: نظرسنجی‌های D3 Systems -- برای دولت آمریکا
% ============================================================
\section{نظرسنجی‌های \lr{D3 Systems} برای وزارت خارجه و ارتش آمریکا}

\subsection{ویژگی خاص}
این نظرسنجی‌ها توسط \lr{D3 Systems} و \lr{KA Research} برای 
وزارت خارجه (\lr{State Department}) و ارتش آمریکا (\lr{DoD}) انجام شدند. 
بخشی از آن‌ها محرمانه بود و بعدها توسط \lr{WikiLeaks} یا 
درخواست‌های \lr{FOIA} منتشر شد.

\begin{table}[H]
\centering
\caption{مشخصات نظرسنجی‌های \lr{D3/KA} برای دولت آمریکا}
\label{tab:d3-method}
\begin{tabularx}{\textwidth}{>{\bfseries}r X}
\toprule
\textbf{ویژگی} & \textbf{شرح} \\
\midrule
مجری & \lr{D3 Systems} (وین، ویرجینیا) و \lr{KA Research} (استانبول) \\
سفارش‌دهنده & وزارت خارجه و وزارت دفاع آمریکا \\
فراوانی & تقریباً هر ۳ ماه (بیش از ۲۰ موج بین ۲۰۰۴ تا ۲۰۱۱) \\
حجم نمونه & \ptho{۲}{۰۰۰}--\ptho{۵}{۰۰۰} \\
ویژگی & پرسش‌های عملیاتی--نظامی، برخی محرمانه \\
\bottomrule
\end{tabularx}
\end{table}

\subsection{یافته‌های کلیدی افشاشده}

\begin{table}[H]
\centering
\caption{یافته‌های کلیدی نظرسنجی‌های محرمانه‌ی \lr{D3/KA}}
\label{tab:d3-findings}
\small
\begin{tabularx}{\textwidth}{>{\bfseries\small}r >{\small}c >{\small}X}
\toprule
\textbf{پرسش} & \textbf{تاریخ} & \textbf{یافته} \\
\midrule
حملات علیه ائتلاف مشروع & ۲۰۰۶ & ۶۱٪ عراقی‌ها حملات علیه نیروهای آمریکایی را تأیید کردند \\
\addlinespace
اعتماد به آمریکا & ۲۰۰۷ & فقط ۱۸٪ به آمریکا اعتماد داشتند (پایین‌ترین سطح) \\
\addlinespace
آینده‌ی بهتر & ۲۰۰۸ & ۵۵٪ معتقد بودند زندگی فرزندانشان بهتر خواهد بود \\
\addlinespace
خروج آمریکا & ۲۰۰۸ & ۷۳٪ خواهان خروج بر اساس جدول زمانی بودند \\
\addlinespace
اعتماد به ارتش عراق & ۲۰۰۹ & ۷۰٪ به ارتش عراق اعتماد داشتند \\
\addlinespace
اعتماد به مالکی & ۲۰۱۰ & ۴۸٪ عملکرد مالکی را مثبت ارزیابی کردند \\
\addlinespace
آمادگی ارتش عراق & ۲۰۱۱ & ۵۵٪ معتقد بودند ارتش عراق آماده‌ی جایگزینی آمریکاست \\
\bottomrule
\end{tabularx}
\end{table}


% ============================================================
% بخش ۱۰: نظرسنجی‌های Zogby International
% ============================================================
\section{نظرسنجی‌های \lr{Zogby International} (۲۰۰۴--۲۰۰۸)}

\subsection{ویژگی خاص}
جیمز زاگبی (\lr{James Zogby})، مؤسس \lr{Arab American Institute}، 
نظرسنجی‌هایی را در عراق و کل جهان عرب انجام داد که به‌ویژه بر 
ادراک‌های فرهنگی و هویتی تمرکز داشت.

\textbf{پرسش کلیدی:}
\begin{quote}
\lr{``Has the Iraq war brought more or less democracy to the Middle East?''}
\end{quote}

\begin{table}[H]
\centering
\caption{«آیا جنگ عراق دموکراسی بیشتری به خاورمیانه آورده؟» -- \lr{Zogby}}
\label{tab:zogby-democracy}
\begin{tabularx}{\textwidth}{>{\bfseries}r *{4}{>{\centering\arraybackslash}X}}
\toprule
\textbf{کشور} & \textbf{بیشتر} & \textbf{کمتر} & \textbf{بی‌تأثیر} & \textbf{سال} \\
\midrule
عراق (کل) & ۲۸٪ & ۴۵٪ & ۲۲٪ & ۲۰۰۸ \\
عراق (شیعه) & ۳۵٪ & ۳۸٪ & ۲۲٪ & ۲۰۰۸ \\
عراق (سنی) & ۱۰٪ & ۶۵٪ & ۲۰٪ & ۲۰۰۸ \\
عراق (کرد) & ۵۵٪ & ۲۰٪ & ۲۰٪ & ۲۰۰۸ \\
\midrule
اردن & ۱۲٪ & ۵۸٪ & ۲۵٪ & ۲۰۰۸ \\
مصر & ۸٪ & ۶۵٪ & ۲۲٪ & ۲۰۰۸ \\
عربستان & ۱۰٪ & ۵۵٪ & ۲۸٪ & ۲۰۰۸ \\
لبنان & ۱۵٪ & ۵۲٪ & ۲۸٪ & ۲۰۰۸ \\
\bottomrule
\end{tabularx}
\end{table}


% ============================================================
% بخش ۱۱: جمع‌بندی تطبیقی نظرسنجی‌ها
% ============================================================
\section{جمع‌بندی تطبیقی: پنج درس از نظرسنجی‌ها}

\subsection{درس اول: هویت فرقه‌ای قوی‌ترین پیش‌بین است}

\begin{figure}[H]
\centering
\begin{tikzpicture}
\begin{axis}[
  ybar,
  width=15cm,
  height=9cm,
  bar width=12pt,
  xlabel={\rl{پرسش}},
  ylabel={\rl{درصد پاسخ مثبت}},
  symbolic x coords={
    {\rl{ارزش سرنگونی}},
    {\rl{جهت درست}},
    {\rl{زندگی بهتر}},
    {\rl{امنیت خوب}},
    {\rl{اعتماد دولت}},
    {\rl{حملات مشروع}}
  },
  xtick=data,
  xticklabel style={rotate=30, anchor=east, font=\small},
  ymin=0, ymax=100,
  legend pos=north east,
  legend style={font=\small},
  grid=major,
  grid style={dashed, gray!25},
  title={\rl{\textbf{شکاف فرقه‌ای در پاسخ‌ها (۲۰۰۹)}}},
  nodes near coords,
  every node near coord/.append style={font=\tiny},
]

\addplot[fill=green!60] coordinates {
  ({\rl{ارزش سرنگونی}},85) 
  ({\rl{جهت درست}},68) 
  ({\rl{زندگی بهتر}},70) 
  ({\rl{امنیت خوب}},85)
  ({\rl{اعتماد دولت}},55)
  ({\rl{حملات مشروع}},3)
};
\addlegendentry{\rl{کردها}}

\addplot[fill=blue!60] coordinates {
  ({\rl{ارزش سرنگونی}},75) 
  ({\rl{جهت درست}},48) 
  ({\rl{زندگی بهتر}},58) 
  ({\rl{امنیت خوب}},55)
  ({\rl{اعتماد دولت}},55)
  ({\rl{حملات مشروع}},22)
};
\addlegendentry{\rl{شیعیان}}

\addplot[fill=red!60] coordinates {
  ({\rl{ارزش سرنگونی}},23) 
  ({\rl{جهت درست}},18) 
  ({\rl{زندگی بهتر}},22) 
  ({\rl{امنیت خوب}},40)
  ({\rl{اعتماد دولت}},15)
  ({\rl{حملات مشروع}},65)
};
\addlegendentry{\rl{سنی‌ها}}

\end{axis}
\end{tikzpicture}
\caption{شکاف فرقه‌ای در پاسخ‌های کلیدی -- ۲۰۰۹}
\label{fig:sectarian-gap-2009}
\end{figure}

\subsection{درس دوم: فرسایش خوش‌بینی اولیه}

\begin{table}[H]
\centering
\caption{فرسایش خوش‌بینی: مقایسه‌ی نقطه‌ی اوج با نقطه‌ی حضیض}
\label{tab:optimism-erosion}
\begin{tabularx}{\textwidth}{>{\bfseries\small}r >{\small}c >{\small}c >{\small}c >{\small}c}
\toprule
\textbf{شاخص} & \textbf{اوج} & \textbf{سال اوج} & \textbf{حضیض} & \textbf{سال حضیض} \\
\midrule
خوش‌بینی به آینده & ۶۷٪ & ۲۰۰۳ & ۱۵٪ & ۲۰۱۹ \\
اعتماد به دولت & ۵۰٪ & ۲۰۰۴ & ۱۰٪ & ۲۰۱۹ \\
جهت کشور درست & ۵۱٪ & ۲۰۰۴ & ۱۵٪ & ۲۰۱۹ \\
رضایت از حمله & ۶۱٪ & ۲۰۰۹ & ۲۲٪ & ۲۰۰۷ \\
اعتماد به احزاب & ۳۵٪ & ۲۰۰۴ & ۵٪ & ۲۰۱۹ \\
\bottomrule
\end{tabularx}
\end{table}

\subsection{درس سوم: پایداری شگفت‌انگیز تقاضا برای دموکراسی}

\infobox{boxgreen}{پارادوکس دموکراسی عراقی}{
یکی از شگفت‌انگیزترین یافته‌ها آن است که علی‌رغم:
\begin{itemize}[rightmargin=0.5cm]
\item سرخوردگی عمیق از عملکرد نظام سیاسی
\item فروپاشی اعتماد به نهادهای حکومتی
\item تجربه‌ی فاجعه‌بار جنگ فرقه‌ای و داعش
\end{itemize}
حمایت عراقی‌ها از \textbf{اصل دموکراسی} هرگز به زیر ۴۸٪ نرفته است (بارومتر عرب). 
عراقی‌ها بین \textbf{ایده‌ی دموکراسی} و \textbf{عملکرد دموکراسی} تمایز قائل‌اند: 
آن‌ها از دومی سرخورده‌اند اما از اولی دست نکشیده‌اند. 

این یافته نشان می‌دهد سرمایه‌ی اجتماعی برای دموکراسی‌سازی هنوز وجود دارد، 
اما اعتمادسازی نهادی ضروری است.
}

\subsection{درس چهارم: شکاف نسلی فزاینده}

\begin{table}[H]
\centering
\caption{شکاف نسلی: جوانان در برابر مسن‌ترها (۲۰۱۹)}
\label{tab:generation-gap}
\begin{tabularx}{\textwidth}{>{\bfseries\small}r >{\centering\arraybackslash\small}X >{\centering\arraybackslash\small}X >{\centering\arraybackslash\small}X}
\toprule
\textbf{شاخص} & \textbf{۱۸--۲۹} & \textbf{۳۰--۴۹} & \textbf{۵۰+} \\
\midrule
اعتماد به دولت & ۷٪ & ۱۰٪ & ۱۵٪ \\
قصد مهاجرت & ۴۵٪ & ۲۸٪ & ۱۰٪ \\
فساد بزرگ‌ترین مشکل & ۵۰٪ & ۴۵٪ & ۳۵٪ \\
دموکراسی بهترین نظام & ۵۲٪ & ۵۸٪ & ۵۵٪ \\
شرکت در اعتراضات & ۳۰٪ & ۱۵٪ & ۵٪ \\
نارضایتی از محاصصه & ۷۵٪ & ۶۵٪ & ۵۰٪ \\
هویت «عراقی» اول & ۵۵٪ & ۴۵٪ & ۳۵٪ \\
\bottomrule
\end{tabularx}
\end{table}

\infobox{boxblue}{تحلیل نسلی}{
نسل پس از ۲۰۰۳ (متولدان ۱۹۹۰ به بعد) مشخصات متفاوتی از نسل‌های قبل دارد:
\begin{itemize}[rightmargin=0.5cm]
\item \textbf{هویت ملی قوی‌تر:} ۵۵٪ هویت «عراقی» را بر هویت فرقه‌ای ترجیح می‌دهند
\item \textbf{سرخوردگی عمیق‌تر:} کمترین اعتماد به نهادها
\item \textbf{تمایل بالای مهاجرت:} ۴۵٪ خواهان مهاجرت هستند
\item \textbf{رادیکالیسم مدنی:} بیشترین مشارکت در اعتراضات (جنبش تشرین)
\item \textbf{رد محاصصه:} ۷۵٪ مخالف نظام سهمیه‌بندی فرقه‌ای
\item این نسل ممکن است موتور تغییرات آینده‌ی عراق باشد
\end{itemize}
}

\subsection{درس پنجم: زنان عراقی -- صداهای ناشنیده}

\begin{table}[H]
\centering
\caption{تفاوت‌های جنسیتی در نظرسنجی‌ها (میانگین ۲۰۰۹--۲۰۱۹)}
\label{tab:gender-gap}
\begin{tabularx}{\textwidth}{>{\bfseries\small}r >{\centering\arraybackslash\small}X >{\centering\arraybackslash\small}X >{\centering\arraybackslash\small}X}
\toprule
\textbf{شاخص} & \textbf{مرد} & \textbf{زن} & \textbf{تفاوت} \\
\midrule
احساس امنیت شبانه & ۵۲٪ & ۳۵٪ & ۱۷ واحد \\
رضایت از وضعیت اقتصادی & ۳۵٪ & ۲۸٪ & ۷ واحد \\
مشارکت در انتخابات & ۶۵٪ & ۵۵٪ & ۱۰ واحد \\
اعتقاد به اینکه رأی تأثیر دارد & ۴۵٪ & ۳۵٪ & ۱۰ واحد \\
موافقت با حقوق برابر زنان & ۳۵٪ & ۵۵٪ & ۲۰ واحد \\
نگرانی از خشونت خانگی & ۲۵٪ & ۵۵٪ & ۳۰ واحد \\
تمایل به مهاجرت & ۳۵٪ & ۲۵٪ & ۱۰ واحد \\
\bottomrule
\end{tabularx}
\end{table}


% ============================================================
% بخش ۱۲: فهرست جامع همه‌ی پرسش‌های کلیدی
% ============================================================
\section{فهرست جامع: تمام پرسش‌های کلیدی نظرسنجی‌ها}

\begin{longtable}{>{\scriptsize\bfseries}p{1.5cm} >{\scriptsize}p{2.5cm} >{\scriptsize}p{5cm} >{\scriptsize}p{2.5cm} >{\scriptsize}p{2.5cm}}
\caption{فهرست جامع پرسش‌های کلیدی نظرسنجی‌ها در عراق} 
\label{tab:all-questions} \\
\toprule
\textbf{کد} & \textbf{مؤسسه} & \textbf{متن پرسش (انگلیسی)} & \textbf{موضوع} & \textbf{سال‌ها} \\
\midrule
\endfirsthead

\multicolumn{5}{r}{\small ادامه از صفحه‌ی قبل} \\
\toprule
\textbf{کد} & \textbf{مؤسسه} & \textbf{متن پرسش} & \textbf{موضوع} & \textbf{سال‌ها} \\
\midrule
\endhead

\midrule
\multicolumn{5}{l}{\small ادامه در صفحه‌ی بعد} \\
\endfoot

\bottomrule
\endlastfoot

\multicolumn{5}{r}{\textbf{الف. ارزیابی حمله و سرنگونی}} \\
\midrule
Q-A1 & ABC/BBC & \lr{Was ousting Saddam worth it?} & ارزش سرنگونی & ۰۴--۰۹ \\
Q-A2 & ABC/BBC & \lr{Right direction or wrong direction?} & جهت کشور & ۰۴--۰۹ \\
Q-A3 & ABC/BBC & \lr{Life compared to before the war?} & مقایسه‌ی زندگی & ۰۴--۰۹ \\
Q-A4 & Pew & \lr{Favorable opinion of the US?} & نگرش به آمریکا & ۰۴--۱۵ \\
Q-A5 & Zogby & \lr{Has war brought more democracy?} & دموکراسی منطقه‌ای & ۰۴--۰۸ \\
\addlinespace

\multicolumn{5}{r}{\textbf{ب. امنیت}} \\
\midrule
Q-B1 & ABC/BBC & \lr{Security in your area?} & امنیت محلی & ۰۴--۰۹ \\
Q-B2 & Gallup & \lr{Feel safe walking at night?} & امنیت شبانه & ۰۶--۱۱ \\
Q-B3 & ABC/BBC & \lr{Attacks on coalition acceptable?} & مشروعیت حملات & ۰۴--۰۹ \\
Q-B4 & Gallup & \lr{Victim of assault/mugging?} & تجربه‌ی خشونت & ۰۷--۱۱ \\
Q-B5 & ABC/BBC & \lr{US presence helps/hurts security?} & تأثیر حضور آمریکا & ۰۴--۰۹ \\
\addlinespace

\multicolumn{5}{r}{\textbf{پ. دموکراسی و حکمرانی}} \\
\midrule
Q-C1 & Arab Bar. & \lr{Democracy better than other systems?} & مطلوبیت دموکراسی & ۰۶--۲۲ \\
Q-C2 & Pew & \lr{Can democracy work here?} & کارایی دموکراسی & ۰۴--۱۲ \\
Q-C3 & WVS & \lr{Strong leader without parliament?} & رهبر قوی & ۰۴، ۱۳ \\
Q-C4 & Arab Bar. & \lr{Trust in government?} & اعتماد به دولت & ۰۶--۲۲ \\
Q-C5 & Arab Bar. & \lr{Trust in parliament?} & اعتماد به پارلمان & ۰۶--۲۲ \\
Q-C6 & ABC/BBC & \lr{Confidence in political parties?} & اعتماد به احزاب & ۰۴--۰۹ \\
Q-C7 & IRI & \lr{Rate PM performance?} & عملکرد نخست‌وزیر & ۰۴--۱۳ \\
Q-C8 & IRI & \lr{Vote makes a difference?} & تأثیر رأی & ۰۵--۱۳ \\
Q-C9 & NDI & \lr{Best government structure?} & ساختار حکومت & ۰۴--۰۹ \\
\addlinespace

\multicolumn{5}{r}{\textbf{ت. اقتصاد و خدمات}} \\
\midrule
Q-D1 & Gallup & \lr{Economic conditions getting better?} & اقتصاد محلی & ۰۶--۱۱ \\
Q-D2 & Gallup & \lr{Satisfied with standard of living?} & رضایت از زندگی & ۰۶--۱۱ \\
Q-D3 & IRI & \lr{Rate electricity availability?} & کیفیت برق & ۰۴--۱۳ \\
Q-D4 & IRI & \lr{Rate water availability?} & کیفیت آب & ۰۴--۱۳ \\
Q-D5 & IRI & \lr{Rate healthcare?} & بهداشت & ۰۴--۱۳ \\
Q-D6 & IRI & \lr{Rate employment opportunities?} & اشتغال & ۰۴--۱۳ \\
Q-D7 & Arab Bar. & \lr{Government performance on corruption?} & مبارزه با فساد & ۱۳--۲۲ \\
\addlinespace

\multicolumn{5}{r}{\textbf{ث. هویت و ارزش‌ها}} \\
\midrule
Q-E1 & ABC/BBC & \lr{Primary identity?} & هویت اول & ۰۷--۰۹ \\
Q-E2 & NDI & \lr{Can groups live together?} & همزیستی & ۰۴--۰۹ \\
Q-E3 & WVS & \lr{Importance of religion?} & اهمیت مذهب & ۰۴، ۱۳ \\
Q-E4 & WVS & \lr{Men better political leaders?} & نگرش جنسیتی & ۰۴، ۱۳ \\
Q-E5 & Pew & \lr{Concern about Sunni-Shia tensions?} & تنش فرقه‌ای & ۰۶--۱۲ \\
Q-E6 & Gallup & \lr{Freedom to choose life?} & آزادی فردی & ۰۶--۱۱ \\
\addlinespace

\multicolumn{5}{r}{\textbf{ج. اولویت‌ها و چالش‌ها}} \\
\midrule
Q-F1 & ABC/BBC & \lr{Biggest problem facing country?} & اولویت اول & ۰۴--۰۹ \\
Q-F2 & Arab Bar. & \lr{Most important challenge?} & مهم‌ترین چالش & ۰۶--۲۲ \\
Q-F3 & Gallup & \lr{Cantril ladder: present?} & رضایت فعلی & ۰۶--۱۱ \\
Q-F4 & Gallup & \lr{Cantril ladder: 5 years?} & امید به آینده & ۰۶--۱۱ \\

\end{longtable}